\chapter{المتجهات والمستقيمات والمستويات في الفضاء}\label{ch:g12:space}

تصير الهندسة الفضائية \omterm{def:g12:seq:arith-geom}{حسابية} متى توافرت المتجهات و
\omterm{def:g10:coordgeom:system}{الإحداثيات}: فتنال المستقيمات تمثيلات وسيطية، وتنال المستويات
معادلات ديكارتية، ويقيس \omterm{def:g12:space:dot}{الجداء السلّمي} الزوايا والمسافات.
ويبني هذا الفصل هذه العدّة ويستعملها لحل مسائل \omterm{def:g10:numbers:interunion}{التقاطع}
والمسافة الكلاسيكية.

\section{المتجهات في الفضاء}

تطيع المتجهات في الفضاء القواعد نفسها التي تطيعها في المستوي: فهي تُجمع وتُضرب في أعداد،
ويكون $\vect{AB} = \vect{CD}$ بالضبط عندما يكون $ABDC$ متوازي أضلاع. وفي
\omterm{def:g10:coordgeom:system}{نظام إحداثيات} $(O; \vec\imath, \vec\jmath, \vec k)$، يكون لكل \omterm{def:g11:vect:vector}{متجهة}
\omterm{def:g10:coordgeom:system}{إحداثيات} $(x, y, z)$.

\begin{definition}[الاستقامية والتوافق في مستو]\label{def:g12:space:collinear}
تكون متجهتان \emph{مستقيميتين} إذا كانت إحداهما مضاعفًا للأخرى. وتكون ثلاث
متجهات $\vec u, \vec v, \vec w$ \emph{متوافقة في مستوٍ}\index{متوافقة في مستو} إذا أمكن كتابة
إحداها تركيبًا من الأخريين، وليكن مثلًا
$\vec w = a\vec u + b\vec v$.
\end{definition}

\begin{proposition}\label{prop:g12:space:basis}
إذا لم تكن $\vec u, \vec v, \vec w$ \omterm{def:g12:space:collinear}{متوافقة في مستوٍ}، فإن كل \omterm{def:g11:vect:vector}{متجهة} في الفضاء
تتفكك على نحو وحيد على الصورة $x\vec u + y\vec v + z\vec w$.
\end{proposition}

\begin{proof}[فكرة البرهان]
عبر طرف \omterm{def:g11:vect:vector}{المتجهة} المراد تفكيكها، ارسم المستقيم الموازي
\omterm{def:g11:vect:vector}{للمتجهة} $\vec w$؛ فهو يلقى مستوي $\vec u, \vec v$ في نقطة واحدة،
تقسم \omterm{def:g11:vect:vector}{المتجهة} إلى مركّبة في ذلك المستوي (وهي وحيدة $x\vec u + y\vec
v$، بهندسة المستوي) ومركّبة على طول $\vec w$. وأما الوحدانية: فإن فرق
تفكيكين سيعبّر عن إحدى المتجهات بدلالة الأخريين،
وهذا يناقض عدم التوافق في مستو.
\end{proof}

\section{المستقيمات والمستويات}

\begin{definition}[التمثيل الوسيطي لمستقيم]\label{def:g12:space:line}
المستقيم المار بالنقطة $A(x_A, y_A, z_A)$ وذو \omterm{def:g11:vect:direction}{متجهة التوجيه}
$\vec u (a, b, c) \neq \vec 0$ هو مجموعة النقط
\[
M(x_A + ta,\; y_A + tb,\; z_A + tc), \qquad t \in \R .
\]
\end{definition}

\begin{definition}[المستوي]\label{def:g12:space:plane}
المستوي المار بالنقطة $A$ والموجَّه بمتجهتين غير \omterm{def:g12:space:collinear}{مستقيميتين}
$\vec u, \vec v$ هو مجموعة النقط $M$ التي تحقق
$\vect{AM} = s\vec u + t\vec v$ حيث $(s, t) \in \R^2$.
\end{definition}

\begin{omfigure}
\begin{tikzpicture}[scale=1.05]
  \draw (-0.5,-0.25) -- (3.4,1.7);
  \draw[-Stealth, omDef, thick] (1,0.5) -- (1.9,0.95)
    node[midway, above left, font=\small] {$\vec u$};
  \fill (0.1,0.05) circle (1.2pt) node[below right, font=\small] {$t=-1$};
  \fill (1,0.5) circle (1.2pt) node[above left, font=\small] {$A$};
  \fill (1.9,0.95) circle (1.2pt) node[below right, font=\small] {$t=1$};
  \fill (2.8,1.4) circle (1.2pt) node[below right, font=\small] {$t=2$};
\end{tikzpicture}
\qquad\quad
\begin{tikzpicture}[scale=1.05]
  \draw (0,0) -- (4,0.5) -- (5.2,1.9) -- (1.2,1.4) -- cycle;
  \node[font=\small] at (4.75,1.6) {$\mathcal P$};
  \coordinate (A) at (1.3,0.55);
  \coordinate (M) at (3.77,1.38);
  \draw[-Stealth, omDef, thick] (A) -- ++(1.31,0.16)
    node[midway, below, font=\small] {$\vec u$};
  \draw[-Stealth, omDef, thick] (A) -- ++(0.42,0.49)
    node[midway, left, font=\small] {$\vec v$};
  \draw[black, dashed, thin] (3.27,0.79) -- (M) -- (1.80,1.14);
  \fill (A) circle (1.2pt) node[below left, font=\small] {$A$};
  \fill (M) circle (1.2pt) node[above, font=\small] {$M$};
\end{tikzpicture}

{\small على اليسار: المستقيم المار بالنقطة $A$ وتوجيهه $\vec u$، مدرَّجًا
\omterm{def:g11:stat:median}{بالوسيط} $t$. وعلى اليمين: المستوي المار بالنقطة $A$ والموجَّه بالمتجهتين $\vec u$ و
$\vec v$؛ وتُبلَغ كل نقطة $M$ من المستوي على الصورة
$\vect{AM} = s\vec u + t\vec v$.}
\end{omfigure}

\begin{proposition}[الأوضاع النسبية]\label{prop:g12:space:positions}
مستويان متمايزان إما متوازيان وإما يتقاطعان في مستقيم. والمستقيم مع
المستوي إما متوازيان (وقد يكون المستقيم محتوى فيه) وإما يتقاطعان في نقطة
واحدة. وأما مستقيمان في الفضاء فقد يكونان متقاطعين، أو متوازيين تمامًا، أو منطبقين،
أو \emph{متخالفين} (أي غير متوافقين في مستوٍ).
\end{proposition}

\begin{proof}[مخطط]
هذه عبارات عن الوقوع؛ وتؤول دراسة كل حالة، \omterm{def:g10:coordgeom:system}{بالإحداثيات}،
إلى حل \omterm{def:g10:lines:system}{جملة خطية} وعدّ حلولها --- انظر
\cref{met:g12:space:intersections}.
\end{proof}

\section{الجداء السلّمي}

\begin{definition}[الجداء السلّمي]\label{def:g12:space:dot}
\emph{الجداء السلّمي}\index{جداء سلّمي} للمتجهتين $\vec u$ و $\vec v$ هو
\[
\vec u \cdot \vec v = \tfrac12\left(\norm{\vec u + \vec v}^2
- \norm{\vec u}^2 - \norm{\vec v}^2\right),
\]
حيث $\norm{\vec u}$ طول $\vec u$. وإذا كانت المتجهتان غير معدومتين،
فإن $\vec u \cdot \vec v = \norm{\vec u}\,\norm{\vec v}\cos\theta$
حيث $\theta$ الزاوية بينهما؛ وفي \omterm{def:g10:coordgeom:system}{نظام إحداثيات}
\emph{متعامد ومتجانس}،
\[
\vec u \cdot \vec v = xx' + yy' + zz' .
\]
وتكون متجهتان \emph{\omterm{def:g11:scal:orthogonal}{متعامدتين}} إذا كان جداؤهما السلّمي $0$.
\end{definition}

\begin{proposition}[الخواص]\label{prop:g12:space:dotprops}
\omterm{def:g12:space:dot}{الجداء السلّمي} متناظر ($\vec u\cdot\vec v = \vec v\cdot\vec u$)،
وخطي ثنائيًا ($(a\vec u + b\vec v)\cdot \vec w = a\,\vec u\cdot\vec w +
b\,\vec v\cdot \vec w$)، و
$\vec u \cdot \vec u = \norm{\vec u}^2 \geq 0$.
\end{proposition}

\begin{proof}
في \omterm{def:g10:coordgeom:system}{نظام متعامد ومتجانس}، تكون الخواص الثلاث كلها مباشرة على الصيغة
$xx' + yy' + zz'$، التي تنتج هي نفسها من الصيغة المعرِّفة ومن
حساب الأطوال بفيثاغورس:
$\norm{\vec u}^2 = x^2 + y^2 + z^2$.
\end{proof}

\begin{definition}[المتجهة الناظمة]\label{def:g12:space:normal}
\emph{المتجهة الناظمة}\index{متجهة ناظمة} لمستوٍ $\mathcal P$ هي
\omterm{def:g11:vect:vector}{متجهة} غير معدومة متعامدة مع كل \omterm{def:g11:vect:vector}{متجهة} تقع في $\mathcal P$ (ويكفي
أن تكون متعامدة مع اتجاهين غير مستقيميين من
$\mathcal P$).
\end{definition}

\begin{omfigure}
\begin{tikzpicture}[scale=1.1]
  \draw (0,0) -- (4,0.5) -- (5.2,1.9) -- (1.2,1.4) -- cycle;
  \node[font=\small] at (4.75,1.6) {$\mathcal P$};
  \coordinate (H) at (2.6,0.95);
  \coordinate (M0) at (2.86,3.03);
  \draw[-Stealth, omThm, thick] (H) -- (2.73,1.99)
    node[midway, right=1pt, font=\small] {$\vec n$};
  \draw[black, dashed, thin] (2.73,1.99) -- (M0);
  \draw[thin] (2.79,0.974) -- (2.815,1.174) -- (2.625,1.15);
  \draw[black, dashed, thin] (H) -- (1.1,0.55);
  \fill (H) circle (1.2pt) node[below right, font=\small] {$H$};
  \fill (M0) circle (1.2pt) node[above, font=\small] {$M_0$};
  \fill (1.1,0.55) circle (1.2pt) node[below left, font=\small] {$A$};
\end{tikzpicture}

{\small \omterm{def:g12:space:normal}{المتجهة الناظمة} $\vec n$ للمستوي $\mathcal P$ متعامدة
مع كل اتجاه من $\mathcal P$. والنقطة $H$، وهي مسقط
العمود من $M_0$، تحقق المسافة من $M_0$ إلى المستوي.}
\end{omfigure}

\begin{theorem}[المعادلة الديكارتية لمستو]\label{thm:g12:space:planeeq}
في \omterm{def:g10:coordgeom:system}{نظام إحداثيات} متعامد ومتجانس، للمستوي المار بالنقطة
$A$ وذي \omterm{def:g12:space:normal}{المتجهة الناظمة} $\vec n(a, b, c)$ \omterm{def:g10:algebra:equation}{المعادلة}
\[
a x + b y + c z + d = 0,
\]
حيث $d = -(ax_A + by_A + cz_A)$. وبالعكس، تعرّف كل \omterm{def:g10:algebra:equation}{معادلة} كهذه مع
$(a,b,c) \neq (0,0,0)$ مستويًا متجهته الناظمة $(a, b, c)$.
\end{theorem}

\begin{proof}
$M \in \mathcal P \iff \vect{AM} \perp \vec n \iff
a(x - x_A) + b(y - y_A) + c(z - z_A) = 0$، وهذا ينشر إلى \omterm{def:g10:algebra:equation}{المعادلة}.
وبالعكس، خذ أي نقطة حل $A$ \omterm{def:g10:algebra:equation}{للمعادلة}؛ فالحساب نفسه
بالمقلوب يبيّن أن مجموعة الحلول هي
$\{M : \vect{AM}\cdot\vec n = 0\}$، أي المستوي المار بالنقطة $A$ والناظم إلى
$\vec n$.
\end{proof}

\begin{method}[التقاطعات عمليًا]\label{met:g12:space:intersections}
\begin{itemize}
  \item \emph{مستقيم $\cap$ مستو}: عوّض \omterm{def:g10:coordgeom:system}{الإحداثيات} الوسيطية
        للمستقيم في \omterm{def:g10:algebra:equation}{معادلة} المستوي؛ وحل بالنسبة إلى $t$ (حل واحد:
        نقطة واحدة؛ ومساواة مستحيلة $0 = c$ مع $c \neq 0$: توازٍ؛ و $0 = 0$: المستقيم محتوى).
  \item \emph{مستو $\cap$ مستو}: حل جملة المعادلتين،
        بالتوسيط بإحداثي حر واحد؛ فالحل مستقيم (أو يكون
        المستويان متوازيين عندما تكون المتجهتان الناظمتان \omterm{def:g12:space:collinear}{مستقيميتين}).
  \item \emph{التحقق من التعامد}: يكون مستقيم $\perp$ مستو إذا وفقط إذا كان توجيهه
        مستقيميًا مع الناظمة؛ ويكون مستويان متعامدين إذا وفقط إذا كانت
        ناظمتاهما \omterm{def:g11:scal:orthogonal}{متعامدتين}.
\end{itemize}
\end{method}

\begin{proposition}[المسافة من نقطة إلى مستو]\label{prop:g12:space:distance}
\index{مسافة!من نقطة إلى مستو}
في \omterm{def:g10:coordgeom:system}{نظام متعامد ومتجانس}، تكون المسافة من $M_0(x_0, y_0, z_0)$ إلى
المستوي $\mathcal P : ax + by + cz + d = 0$ هي
\[
\operatorname{dist}(M_0, \mathcal P)
= \frac{\abs{ax_0 + by_0 + cz_0 + d}}{\sqrt{a^2 + b^2 + c^2}}.
\]
\end{proposition}

\begin{proof}
ليكن $H$ المسقط العمودي للنقطة $M_0$ على $\mathcal P$: أي مسقط
العمود، فتكون $\vect{HM_0}$ مستقيمية مع الناظمة الواحدية
$\frac{\vec n}{\norm{\vec n}}$، وتكون المسافة
$\abs{\vect{HM_0}\cdot \frac{\vec n}{\norm{\vec n}}}$. ومن أجل أي
$H \in \mathcal P$،
$\vect{HM_0}\cdot\vec n = ax_0 + by_0 + cz_0 - (ax_H + by_H + cz_H)
= ax_0 + by_0 + cz_0 + d$ (باستعمال \omterm{def:g10:algebra:equation}{المعادلة} عند $H$)، ومنه الصيغة
بعد القسمة على $\norm{\vec n} = \sqrt{a^2+b^2+c^2}$.
\end{proof}

\begin{example}\label{ex:g12:space:distance}
المسافة من المبدأ إلى المستوي $x + 2y + 2z - 6 = 0$:
$\frac{\abs{-6}}{\sqrt{1 + 4 + 4}} = \frac63 = 2$.
\end{example}

\section{تمارين}

في كل التمارين يكون \omterm{def:g10:coordgeom:system}{نظام الإحداثيات} \omterm{def:g10:coordgeom:system}{متعامدًا ومتجانسًا}.

\begin{exercise}[$\star$]\label{exo:g12:space:1}
لتكن $A(1, 0, 2)$ و $B(3, 1, 1)$ و $C(2, -1, 3)$. فهل النقط $A$ و $B$ و
$C$ على استقامة واحدة؟ أعطِ تمثيلًا وسيطيًا للمستقيم $(AB)$.
\end{exercise}

\begin{exercise}[$\star$]\label{exo:g12:space:2}
حدّد \omterm{def:g10:algebra:equation}{معادلة} ديكارتية للمستوي المار بالنقطة $A(1, 1, 0)$ وذي
\omterm{def:g12:space:normal}{المتجهة الناظمة} $\vec n(2, -1, 3)$. فهل تنتمي النقطة $B(0, 2, 1)$ إليه؟
\end{exercise}

\begin{exercise}[$\star$]\label{exo:g12:space:3}
احسب الزاوية عند $A$ (مقرَّبةً إلى أقرب درجة) في المثلث ذي
الرؤوس $A(0,0,0)$ و $B(1, 1, 0)$ و $C(1, 0, 1)$.
\end{exercise}

\begin{exercise}[$\star\star$]\label{exo:g12:space:4}
اعتبر المستقيم $\mathcal D$ المار بالنقطة $A(1, 2, 0)$ وذا التوجيه
$\vec u(1, -1, 2)$، والمستوي $\mathcal P : 2x + y - z + 1 = 0$.
حدّد $\mathcal D \cap \mathcal P$.
\end{exercise}

\begin{exercise}[$\star\star$]\label{exo:g12:space:5}
ليكن $\mathcal P : x - y + 2z = 1$ و $\mathcal Q : 2x + y + z = 4$.
بيّن أن $\mathcal P$ و $\mathcal Q$ يتقاطعان في مستقيم وأعطِ
تمثيلًا وسيطيًا له.
\end{exercise}

\begin{exercise}[$\star\star$]\label{exo:g12:space:6}
بيّن أن المستقيمين
\[
\mathcal D_1 : (x, y, z) = (1 + t,\; 2t,\; 3 - t), \qquad
\mathcal D_2 : (x, y, z) = (2 + s,\; 1 + s,\; 1 + 2s)
\]
متخالفان (أي غير متوافقين في مستوٍ).
\end{exercise}

\begin{exercise}[$\star\star$]\label{exo:g12:space:7}
ليكن $\mathcal P : 2x - y + 2z - 3 = 0$.
\begin{enumerate}
  \item احسب المسافة من $M_0(3, 1, 1)$ إلى $\mathcal P$.
  \item حدّد \omterm{def:g10:coordgeom:system}{إحداثيات} المسقط العمودي $H$ للنقطة
        $M_0$ على $\mathcal P$، وتحقق من المسافة $M_0H$.
\end{enumerate}
\end{exercise}

\begin{exercise}[$\star\star\star$]\label{exo:g12:space:8}
المكعب $ABCDEFGH$ طول ضلعه $1$؛ ضع \omterm{def:g10:coordgeom:system}{الإحداثيات} بحيث
$A(0,0,0)$ و $B(1,0,0)$ و $D(0,1,0)$ و $E(0,0,1)$ (مع
$C = B + D$ و $F = B + E$ و $G = B + D + E$ و $H = D + E$ بوصفها متجهات انطلاقًا من
$A$).
\begin{enumerate}
  \item بيّن أن القطر $(AG)$ عمودي على المستوي $(BDE)$.
  \item احسب المسافة من $A$ إلى المستوي $(BDE)$، والنقطة
        التي يخترقه عندها $(AG)$.
\end{enumerate}
\end{exercise}

\begin{exercise}[$\star\star\star$]\label{exo:g12:space:9}
\emph{(حجم رباعي وجوه.)} لتكن $A(1,1,1)$ و $B(2,1,0)$ و $C(0,2,1)$
و $D(2,2,2)$.
\begin{enumerate}
  \item تحقق من أن $\vect{AB}$ و $\vect{AC}$ و $\vect{AD}$ غير \omterm{def:g12:space:collinear}{متوافقة في مستوٍ}
        (فالنقط تكوّن رباعي وجوه حقيقيًا).
  \item أوجد \omterm{def:g11:vect:vector}{متجهة} $\vec n$ متعامدة مع $\vect{BC}$ ومع
        $\vect{BD}$ معًا، واستنتج \omterm{def:g10:algebra:equation}{معادلة} ديكارتية للمستوي $(BCD)$.
  \item احسب المسافة من $A$ إلى المستوي $(BCD)$ ومساحة
        المثلث $BCD$، باستعمال الصيغة
        $\text{المساحة} = \frac12\sqrt{\norm{\vec u}^2\norm{\vec v}^2 -
        (\vec u \cdot \vec v)^2}$ من أجل المثلث الذي تولّده
        $\vec u, \vec v$.
  \item استنتج حجم رباعي الوجوه
        ($V = \frac13 \times \text{مساحة القاعدة} \times \text{الارتفاع}$).
\end{enumerate}
\end{exercise}

\section{مسألة: تشريح المكعب}

\begin{problem}[{مسألة نهاية الأسبوع --- مستوٍ يقطع مكعبًا إلى
مسدَّس منتظم، ورباعي وجوه منتظم مختبئ في الزوايا، و
زاوية القطر تشكّل كل ألماسة}]
\label{pb:g12:space:1}
اقطع مكعبًا وأنت تتوقع مربعات ومستطيلات --- ومع ذلك فإن قطعًا
شهيرًا واحدًا ينتج \emph{مسدَّسًا منتظمًا}، وأربعة قطوع بارعة
في الزوايا تترك وراءها \emph{رباعي وجوه منتظمًا}. وكلتا
المفاجأتين، والزاوية $109.5^\circ$ التي تعبدها ذرات الكربون،
تسقطان أمام عدّة هذا الفصل: المستقيمات الوسيطية، ومعادلات المستويات،
\omterm{def:g12:space:dot}{والجداء السلّمي} (\cref{thm:g12:space:planeeq} و
\cref{prop:g12:space:distance}). واستعمل في كل ما يلي مكعب
\cref{exo:g12:space:8}: طول ضلعه $1$، و $A(0,0,0)$ و $B(1,0,0)$
و $D(0,1,0)$ و $E(0,0,1)$ و $C(1,1,0)$ و $F(1,0,1)$ و $H(0,1,1)$
و $G(1,1,1)$.

\medskip
\noindent\textbf{الجزء الأول --- التمكّن.}
\begin{enumerate}
  \item اكتب \omterm{def:g10:algebra:equation}{معادلة} المستوي المار بالنقطة $B(1,0,0)$ وذي
        \omterm{def:g12:space:normal}{المتجهة الناظمة} $(1,1,1)$، وتمثيلًا وسيطيًا
        للمستقيم المار بالنقطة $A$ والموجَّه \omterm{def:g11:vect:vector}{بالمتجهة} $(1,1,1)$، ثم
        نقطة تقاطعهما.
  \item احسب المسافة من $A$ إلى ذلك المستوي
        (\cref{prop:g12:space:distance}).
  \item هل المتجهتان $(1,-1,0)$ و $(1,1,-2)$ متعامدتان؟
  \item ما وضع المستويين $x + y + z = 1$ و
        $2x + 2y + 2z = 5$؟ وما وضع المستقيم المار بالنقطة $A$ والموجَّه \omterm{def:g11:vect:vector}{بالمتجهة}
        $(1,1,0)$ بالنسبة إلى المستوي الأول؟
  \item اسرد منتصفات الحواف الست
        $[BC], [CD], [DH], [HE], [EF], [FB]$ للمكعب، و
        تحقق من أن الستة كلها تحقق \omterm{def:g10:algebra:equation}{المعادلة}
        $x + y + z = \frac32$.
\end{enumerate}

\medskip
\noindent\textbf{الجزء الثاني --- المقطع المسدَّس.}
ليكن $P$ المستوي $x + y + z = \frac32$.
\begin{enumerate}[resume]
  \item بيّن أن $P$ عمودي على القطر الكبير
        $(AG)$ ويمر بمركز المكعب.
  \item احسب المسافات بين منتصفات السؤال 5 المتعاقبة
        حول المقطع: فماذا تجد؟
  \item احسب زاوية مضلع المقطع عند أحد
        رؤوسه (بمتجهتي حافتين متعاقبتين \omterm{def:g12:space:dot}{وجداء
        سلّمي}). واختم: المقطع \emph{مسدَّس
        منتظم}.
  \item احسب محيط المسدَّس ومساحته، وقارن
        المساحة بمساحة وجه من المكعب. (وللسجل: أكبر مقطع
        مستوٍ لمكعب الوحدة هو المستطيل القطري $1 \times \sqrt2$
        الذي مساحته $\sqrt2$ --- فينال المسدَّس المنتظم الفضة.)
  \item أين يخترق القطر $(AG)$ المستوي $P$؟ وتحقق من أن
        نقطة الاختراق هي \omterm{prop:g10:coordgeom:midpoint}{منتصف} $[AG]$.
  \item حرّك القطع: صِف المقاطع
        $x + y + z = c$ لما ينمو $c$ من $0$ إلى $3$ ---
        واحسب المقطع من أجل $c = \frac12$ (أي مضلع،
        وأي حجم؟)، واحكِ التحول الذي يمر
        بمسدَّسك عند $c = \frac32$.
  \item لماذا لا يستطيع أي مستوٍ أن يقطع مكعبًا في مضلع ذي
        \emph{سبعة} أضلاع؟ (أين يجب أن يقع كل ضلع من
        المقطع؟)
\end{enumerate}

\medskip
\noindent\textbf{الجزء الثالث --- رباعي الوجوه المختبئ.}
\begin{enumerate}[resume]
  \item احسب حجم رباعي وجوه الزاوية $ABDE$
        (بالقاعدة والارتفاع وصيغة الثلث التي قبلها الكتاب
        السابق --- و\cref{exo:g12:space:9}
        يحسب مثل هذه الأحجام عمومًا).
  \item الرؤوس الأربعة $B$ و $D$ و $E$ و $G$: احسب المسافات
        الست بين كل زوجين واختم بأن $BDEG$ رباعي وجوه
        \emph{منتظم}. ثم احصل على حجمه بطرح
        رباعيات وجوه الزوايا من المكعب.
  \item احسب المسافة من مركز المكعب إلى
        المستوي $(BDE)$.
  \item زاوية الألماس: احسب الزاوية بين
        القطرين $\vect{AG}$ و $\vect{BH}$ في المكعب. ومتممتها، ونحوها $109.5^\circ$، هي الزاوية بين
        الروابط في الألماس والميثان --- فلماذا تختار أربعة
        أزواج إلكترونية حول ذرة كربون زوايا
        رباعي وجوه السؤال 14؟
\end{enumerate}

\medskip
\noindent\textbf{الجزء الرابع --- ثلاثي الأبعاد حقًا.}
\begin{enumerate}[resume]
  \item أوجد أين يخترق المستقيم المار بالنقطة $B(1,0,0)$ والموجَّه \omterm{def:g11:vect:vector}{بالمتجهة}
        $(0,1,1)$ مستوي المسدَّس $P$.
  \item بيّن أن المستقيمين $(AB)$ و $(EG)$ لا يلتقيان ولا
        يتوازيان: فهما مستقيمان \emph{متخالفان}، وهي ظاهرة
        لا يمكن أن تحدث في مستوٍ.
  \item احسب المسقط العمودي للنقطة $B$ على $P$
        (بالسير من $B$ على طول الناظمة حتى المستوي)، ثم
        تحقق من أن نقطتك تحقق \omterm{def:g10:algebra:equation}{معادلة} $P$.
  \item الخاتمة --- عدّة مساح الفضاء: المستقيمات الوسيطية
        للاختراق، والمتجهات الناظمة للزوايا والمسافات،
        ومعادلات المستويات لنحت المقاطع؛ والمكعب ملعبًا
        أنتجت فيه مسدَّسًا، ورباعي وجوه منتظمًا،
        وزاوية الألماس، وزوجًا من المستقيمات المتخالفة.
        جملة واحدة لكل بند.
\end{enumerate}
\end{problem}
